%% =====================================================================
%%  Bilan 3 — En ville (unités 9 à 12)
%%  Pistes : B3.01 (écoute notée) · B3.02 (expression orale)
%% =====================================================================
\bilan{3}

\noindent{\fontsize{11}{14.6}\selectfont Troisième partie terminée ! Vous savez maintenant demander
votre chemin, faire des achats, commander au café et parler de vos talents. Révisez l'essentiel, puis
faites les huit exercices notés sur 20. Les corrigés sont page~\pageref{cor-b3-1}.\par}

%% ---------------------------------------------------------------------
\rubrique{grammaire}{Révisez}

\begin{paire}
\begin{regle}[L'impératif et le chemin · unité 9]
\en{Turn left.} · \en{Go straight ahead.} · \en{Cross the street.}\par
Négation : \en{Don't take the bus.}\par\vspace{1mm}
\en{Excuse me, where's the bank?} · \en{How do I get to…?} · \en{Is it far?}\par\vspace{1mm}
\en{next to} · \en{opposite} · \en{between} · \en{on the corner of}
\end{regle}
\begin{regle}[Quantités et prix · unité 10]
\en{How much is it? — It's twenty dollars.}\par
\en{How much} + indénombrable · \en{How many} + pluriel\par\vspace{1mm}
\en{some} (affirmation) · \en{any} (négation, question)\par\vspace{1mm}
{\fontsize{9.6}{12}\selectfont \en{thirteen} \phon{θɝːˈtiːn} ≠ \en{thirty} \phon{ˈθɝːti}\par}
\end{regle}
\end{paire}

\begin{paire}
\begin{regle}[Commander · unité 11]
\en{Can I have a coffee, please?} · \en{I'd like a salad.} · \en{I'll have the burger.}\par\vspace{1mm}
\en{Would you like something to drink?}\par
\en{Can we have the check, please?}
\end{regle}
\begin{regle}[\en{can} · unité 12]
\en{I can swim.} · \en{She can't ski.}\par
\en{Can you drive? — Yes, I can. / No, I can't.}\par\vspace{1mm}
\en{very well} · \en{quite well} · \en{not very well} · \en{not at all}\par\vspace{1mm}
{\fontsize{9.6}{12}\selectfont Ni \en{-s}, ni \en{to}, ni \en{do} : \en{she can drive}.\par}
\end{regle}
\end{paire}

%% ---------------------------------------------------------------------
\rubrique{pratique}{Testez-vous}

\exercice[ecrire,niveau=1,points=2]{Le chemin}{Traduisez ces indications. \textit{(0,5 point par
réponse)}}
\begin{questions}
\item Tournez à gauche. \hfill \enb{\trou[60mm]{Turn left.}}
\item Allez tout droit. \hfill \enb{\trou[60mm]{Go straight ahead.}}
\item Traversez la rue. \hfill \enb{\trou[60mm]{Cross the street.}}
\item Ne prenez pas le bus. \hfill \enb{\trou[60mm]{Don't take the bus.}}
\end{questions}

\exercice[ecrire,niveau=2,points=3]{Quantités}{Complétez avec \en{How much}, \en{How many},
\en{some} ou \en{any}. \textit{(0,5 point par réponse)}}
\begin{anglais}
\begin{questions}
\item \trou[24mm]{How much} is this jacket?
\item \trou[24mm]{How many} apples do you need?
\item I'd like \trou[24mm]{some} water, please.
\item There aren't \trou[24mm]{any} eggs in the fridge.
\item Are there \trou[24mm]{any} stores near here?
\item \trou[24mm]{How much} money do you have?
\end{questions}
\end{anglais}

\exercice[ecrire,niveau=1,points=2]{Les prix}{Écrivez chaque prix en toutes lettres.
\textit{(0,5 point par réponse)}}
\begin{questions}[2]
\item \enb{\$15} → \enb{\trou[44mm]{fifteen dollars}}
\item \enb{\$50} → \enb{\trou[44mm]{fifty dollars}}
\item \enb{\$13.99} → \enb{\trou[44mm]{thirteen ninety-nine}}
\item \enb{\$80} → \enb{\trou[44mm]{eighty dollars}}
\end{questions}

\exercice[ecrire,niveau=2,points=3]{Au restaurant}{Complétez le dialogue. \textit{(0,5 point par
réponse)}}
\begin{anglais}
\noindent\begin{tabular}{@{}l@{\hspace{6mm}}l@{}}
\otNomDit{serveur} & Are you ready to \trou{order}?\\[1.8mm]
\otNomDit{lea}     & Yes. \trou{Can} I have the chicken salad, please?\\[1.8mm]
\otNomDit{serveur} & Sure. \trou{Would} you like something to drink?\\[1.8mm]
\otNomDit{lea}     & I'd \trou{like} an orange juice, please.\\[1.8mm]
\otNomDit{serveur} & Here you \trou{are}. Enjoy!\\[1.8mm]
\otNomDit{lea}     & Thank you! \ldots{} Can we have the \trou{check}, please?\\
\end{tabular}
\end{anglais}

\exercice[lire,ecrire,niveau=1,points=2]{Qui sait faire quoi ?}{Regardez les symboles et complétez avec
\en{can} ou \en{can't}. \textit{(0,5 point par réponse)}}
\begin{anglais}
\begin{questions}[2]
\item \otCheck\quad Sam \trou[16mm]{can} play the guitar.
\item \otCroix\quad Jack \trou[16mm]{can't} ski.
\item \otCroix\quad Maya \trou[16mm]{can't} drive.
\item \otCheck\quad Léa \trou[16mm]{can} swim very well.
\end{questions}
\end{anglais}

\exercice[ecouter,niveau=2,points=3,piste=B3.01]{Au magasin}{Léa achète un vêtement. Écoutez, puis
indiquez si chaque phrase est vraie (V) ou fausse (F). \textit{(0,5 point par réponse)}}
\begin{anglais}
\begin{questions}[2]
\item Léa needs a sweater. \hfill\trou[9mm]{V}
\item She's a size small. \hfill\trou[9mm]{F}
\item The sweater is blue. \hfill\trou[9mm]{V}
\item It's forty-nine dollars. \hfill\trou[9mm]{F}
\item She tries it on. \hfill\trou[9mm]{V}
\item She pays cash. \hfill\trou[9mm]{F}
\end{questions}
\end{anglais}

\exercice[lire,niveau=2,points=3]{Le menu}{Lisez le menu, puis indiquez si chaque phrase est vraie (V)
ou fausse (F). \textit{(0,5 point par réponse)}}
\begin{anglais}
\begin{tcolorbox}[enhanced,width=124mm,center,colback=otOrangeTint!40,colframe=otOrangeDark!70,boxrule=0.7pt,
  arc=2mm,left=6mm,right=6mm,top=3.5mm,bottom=3.5mm,fontupper=\fontsize{10.4}{13.4}\selectfont]
\begin{center}{\otBold\fontsize{15}{18}\selectfont\color{otInk}Brooklyn Corner Café}\par\vspace{0.6mm}
{\fontsize{8.4}{10}\selectfont\color{otMuted}Open every day · 7 a.m. – 9 p.m.}\end{center}\vspace{1mm}
\begin{minipage}[t]{0.47\linewidth}
{\otSemi\fontsize{8.6}{10}\selectfont\color{otOrangeDark}\otCaps[12]{Drinks}}\par\vspace{0.8mm}
Coffee\dotfill \$2.50\par Latte\dotfill \$4.75\par Tea\dotfill \$2.25\par Orange juice\dotfill \$3.95\par
\vspace{2mm}{\otSemi\fontsize{8.6}{10}\selectfont\color{otOrangeDark}\otCaps[12]{Desserts}}\par\vspace{0.8mm}
Ice cream\dotfill \$4.50\par Muffin\dotfill \$3.25
\end{minipage}\hfill
\begin{minipage}[t]{0.47\linewidth}
{\otSemi\fontsize{8.6}{10}\selectfont\color{otOrangeDark}\otCaps[12]{Food}}\par\vspace{0.8mm}
Soup of the day\dotfill \$6.50\par Chicken salad\dotfill \$8.95\par Burger and fries\dotfill \$12.50\par
Steak\dotfill \$24.00
\end{minipage}
\par\vspace{2mm}\begin{center}{\fontsize{8.4}{10}\selectfont\color{otMuted}Tax and tip not included}\end{center}
\end{tcolorbox}
\begin{questions}[2]
\item A tea is \$2.25. \hfill\trou[9mm]{V}
\item The burger comes with fries. \hfill\trou[9mm]{V}
\item The steak is \$14. \hfill\trou[9mm]{F}
\item There's a soup on the menu. \hfill\trou[9mm]{V}
\item There are three desserts. \hfill\trou[9mm]{F}
\item The tip is included. \hfill\trou[9mm]{F}
\end{questions}
\end{anglais}

\exercice[ecrire,niveau=3,points=2]{Vos talents}{En quatre phrases, dites ce que vous savez faire et ce
que vous ne savez pas faire, en précisant votre niveau. \textit{(0,5 point par phrase correcte)}}
\lignes{4}

%% ---------------------------------------------------------------------
\rubrique{ecouter}{Parlez}

\exercice[ecouter,parler,niveau=3,piste=B3.02]{Au café}{Vous êtes au comptoir d'un café de New York.
Répondez au serveur \textbf{à voix haute} pendant chaque pause. Cet exercice n'est pas noté : comparez
vos réponses avec les modèles du corrigé.}
\begin{center}\otTableau
\begin{tabular}{@{}>{\otSemi\fontsize{8.6}{10}\selectfont\color{otNavy}}l@{\hspace{3mm}}>{\enOn}p{70mm}<{\enOff}@{\hspace{6mm}}p{78mm}@{}}
 & \otTH{Le serveur dit\ldots} & \otTH{Vous\ldots}\\\midrule
1. & Hi! What can I get you? & commandez une boisson.\\
2. & Would you like something to eat? & commandez quelque chose à manger.\\
3. & For here or to go? & répondez.\\
4. & That's nine dollars. Cash or card? & demandez si vous pouvez payer par carte.\\
5. & Sure! Have a nice day! & remerciez et dites au revoir.\\\bottomrule
\end{tabular}\end{center}

\otScore{20}
